\documentclass[../main.tex]{subfiles}
\begin{document}

\section{Overview}
The previous chapter validated the summary graph empirically; this chapter steps back to
discuss what the findings mean, where the method falls short, and what they amount to as a
whole. It first weighs the results and the method's limitations, then summarizes the
thesis's contributions and outcomes and points to future work.

\section{Discussion and Limitations}
The central finding of Chapter~5 is that the summarization pipeline produces causally
faithful summary graphs: across the analogy graphs, steering a supernode moves its
downstream neighbours in the direction and roughly the magnitude its summary edge predicts,
for both the exact integer-program clusterer and the $K$-means baseline. Faithfulness is
therefore a property of the abstraction rather than of one particular clusterer, which is
the outcome the objective was designed to produce. The comparison between the two clusterers
is more nuanced than a single winner: at a fixed supernode count the simpler baseline is
more faithful per edge, while the integer program exposes more of the circuit's causal
wiring through its causal-preservation term. The two are different operating points on one
trade-off, not a ranking, and the steering metric, which grades a method only on the edges
it chooses to expose, is partly biased toward the method that exposes less structure.

Several limitations bound these conclusions. The method summarizes one prompt at a time and
inherits the scope of the underlying attribution graph, so it says nothing about mechanisms
shared across prompts or about circuits that run through attention rather than the MLP
features the replacement model captures. The exact clusterer solves an integer program whose
size grows with the number of feature pairs and is tractable only on the small pruned graphs
studied here; larger graphs fall back to the approximate spectral and agglomerative
clusterers. The trade-off weight between atomicity and causal preservation was fixed at a
single value in the evaluation, so the operating point at which the per-edge and structural
effects balance has not yet been located; a sweep over this weight is the natural next
experiment. Finally, the evaluation rests on the steering intervention alone: the
complementary LLM-versus-human relevance comparison is still in progress, and feature
meanings are taken from existing labels rather than verified, so the readability of the
summaries has not been measured against human judgment.

\section{Conclusion}
This thesis formalized the problem of summarizing an attribution graph and proposed an
automated three-stage pipeline --- pruning, clustering, and interpreting --- that replaces
the manual reduction of the original circuit-tracing methodology. The problem was stated
through four desiderata and a matching objective, the pipeline was built to optimize that
objective with an exact integer program and scalable alternatives, and the resulting summary
graphs were shown to be causally faithful under steering interventions on the BATS analogy
and multi-hop datasets. Together these results turn a slow, expert-driven analysis into an
objective-driven procedure that runs without human intervention while keeping the summary
faithful to the model. Future work follows from the limitations above: sweeping the
atomicity--causal trade-off to characterize the full Pareto front, scaling the exact solver
or tightening the approximate clusterers for larger graphs, completing the human-judgment
evaluation of readability, and extending the summary beyond a single prompt toward mechanisms
that recur across many.

\section{Chapter Summary}
This chapter discussed the thesis's results and their limits and drew its conclusions. The
summarization pipeline was shown to deliver faithful, automated summary graphs, with the
choice of clusterer setting a trade-off between per-edge faithfulness and exposed structure
rather than a strict ranking. The main limitations --- the single-prompt scope, the
tractability ceiling of the exact solver, the unswept trade-off weight, and the
still-incomplete readability evaluation --- define the agenda for future work, completing
the account of the method begun in Chapter~3 and validated in Chapter~5.

\end{document}
